% Generated from locale/tr/content/lambda-calculus/syntax/terms.tex; do not hand-edit.


\olfileid[tr]{lam}{syn}{trm}
\olsection{Terimler}

Lambda hesabının terimleri, sonsuz bir $\Obj{v_0}$, $\Obj{v_1}$, \dots{}
değişken dağarcığından, ``$\lambd$'' simgesinden ve parantezlerden
tümevarımlı olarak kurulur. Değişkenleri $x$, $y$, $z$, \dots{}; terimleri
ise $M$, $N$, $P$, \dots{} ile göstereceğiz.

\begin{defn}[Terimler] \ollabel{defn:term}
Lambda hesabının \emph{terimler} kümesi tümevarımlı olarak şöyle tanımlanır:
\begin{enumerate}
  \item \ollabel{defn:term-var} $x$ bir değişkense $x$ bir terimdir.
  \item \ollabel{defn:term-abs} $x$ bir değişken ve $M$ bir terimse
    $(\lambd[x][M])$ bir terimdir.
  \item \ollabel{defn:term-app} Hem $M$ hem de $N$ birer terimse $(MN)$ bir
    terimdir.
\end{enumerate}
\end{defn}

Bir $(\lambd[x][M])$ terimi \olref{defn:term-abs} uyarınca oluşturulmuşsa
bunun bir \emph{soyutlama} sonucunda elde edildiğini söyleriz; $\lambd[x]$
içindeki $x$'e de \emph{!!{parameter}} denir. \olref{defn:term-app} uyarınca
oluşturulan bir $(MN)$ terimi ise bir \emph{uygulama} sonucudur.

Yukarıda tanımlanan terimler bütünüyle parantezlenmiştir. Bu gösterim,
$(\lambd[x][((\lambd[x][x])(\lambd[x][(xx)]))])$ teriminin gösterdiği gibi
oldukça hantal olabilir. Parantezlerden kaçınmak için bazı uzlaşımlar
getireceğiz. Bununla birlikte resmî tanım, bir terimin \olref{defn:term}
uyarınca nasıl kurulduğunu belirlemeyi kolaylaştırır. Örneğin
$(\lambd[x][((\lambd[x][x])(\lambd[x][(xx)]))])$ terimini oluşturmanın son
adımı, !!{parameter} $x$ olan bir soyutlama olmak zorundadır. Bu terim,
iki terimin uygulaması olan
$((\lambd[x][x])(\lambd[x][(xx)]))$ teriminden soyutlamayla elde edilir.
Bu iki terimin her biri de birer soyutlama sonucudur ve çözümleme bu biçimde
sürer.

\begin{prob}
$(\lambd[g][(\lambd[x][(g (x x))])
  (\lambd[x][(g (x x))])])$ teriminin nasıl oluşturulduğunu açıklayın.
\end{prob}

